% ============================================================================
% dissolution-theorem.tex  (v3, 2026-08-15)
%
% The dissolution law of the introduction, proved:  under GRH(chi) and
% L(1/2,chi) != 0 for a real primitive character chi mod q, the density
% delta_q(m) := P(X_{m,chi} > 0) of the limiting variable satisfies, as
% m -> infinity,
%     delta_q(m) - 1/2 = (2 pi sigma_m^2)^{-1/2} (1 + O_q(1/(M log M))),
% with a proved second-order refinement, and (for q = 4, under LI, via
% Theorem thm:interp(ii)) the same asymptotics for the logarithmic density
% of the weighted race itself.  This is the theorem forward-referenced as
% \ref{thm:dissolution} in thm-full-range.tex.
%
% Intended placement: \input into the v3 main file AFTER the theorem section
% (as extended by thm-full-range.tex) and AFTER the variance section
% (variance-identity.tex).  It uses the following labels of those files:
%   thm:interp         (theorem, range (-1/2,infty) per thm-full-range.tex)
%   prop:varid         (variance identity, q=4)
%   prop:varasymp      (variance asymptotic, general real primitive chi;
%                       its eq:varidq is the general-q variance identity)
%   rem:tails          (exact tail variances for the numerics)
%   tab:delta          (density table of the introduction)
%   tab:enclose        (certified enclosures; v3/certified-numerics.tex --
%                       referenced in one sentence of the numerics remark;
%                       if that section is not included, drop the sentence)
% Bibliography keys used: Dav00, ANS14, FM13 (all in the main file)
% plus ONE NEW key that must be added to the main bibliography:
%   \bibitem{Fel71} W. Feller, \emph{An Introduction to Probability Theory
%   and Its Applications}, Vol.~II, 2nd ed., Wiley, New York, 1971.
%
% Every numeric in this section (the Bessel constants of Lemma
% \ref{lem:J0facts}, the closed-form values of sigma_m^2 and S_4(m), the
% reference densities, and the entries of Table \ref{tab:dissolution}) was
% recomputed for this draft by
%   paper/v3/verify_dissolution.py
% (mpmath 1.4.1 at 50 working digits with 90 guard digits around the
% Hurwitz-zeta pole at s=1, numpy/scipy for the deterministic Gil--Pelaez
% reference densities; see the script header for the checklist).
%
% Suggested edits to other files once this section lands (not applied here):
%  * weighted-races.tex, Results item 3 ("Dissolution law"): after the
%    displayed law add "now a theorem: see Section~\ref{sec:dissolution}".
%  * weighted-races.tex, "Assumptions and limitations": the clause "the
%    dissolution law is derived at the level of rigor of \cite{FM13}'s
%    heuristic term plus numerical confirmation, and a fully rigorous proof
%    (with error term) is left open" is superseded by
%    Theorem~\ref{thm:dissolution} and should be deleted or rewritten.
%  * variance-identity.tex, Remark rem:varasymp-numerics part (3) ("What is
%    not proved here ... retains the status it had") is superseded by
%    Theorem~\ref{thm:dissolution} and should point to it.
%  * thm-full-range.tex, Theorem thm:interp(iv) final sentence "which this
%    theorem does not prove" correctly anticipates this section; optionally
%    append "(see Theorem~\ref{thm:dissolution})".
% ============================================================================

\section{The dissolution law: from heuristic to theorem}
\label{sec:dissolution}

The introduction stated the dissolution law
$\delta_q(m)-\tfrac12\sim(2\pi\,\mathrm{Var}_q(m))^{-1/2}$ as a
Gaussian-comparison heuristic with numerical confirmation, and
Proposition~\ref{prop:varasymp} made its variance input
$\mathrm{Var}_q(m)=2M\log\tfrac{qM}{2\pi}+O(1)$ rigorous. This section
proves the density half, with a rate, and in fact with a proved
second-order term. The method is direct Fourier inversion: the limiting
variable is $1+V$ with $V$ a sum of independent scaled cosines, its
characteristic function is a product of Bessel factors
$J_0(2a_\gamma t)$, and
\[
\delta_q(m)-\tfrac12\;=\;\frac1\pi\int_0^\infty
\Phi_m(t)\,\frac{\sin t}{t}\,dt,
\qquad
\Phi_m(t)=\prod_{\gamma>0}J_0\bigl(2a_\gamma(m)\,t\bigr).
\]
Near $t=0$ the product is Gaussian to high accuracy, with variance
$\sigma_m^2$ and a fourth-cumulant correction governed by
$S_4(m):=\sum_{\gamma>0}(2a_\gamma(m))^4$; away from $t=0$ the
$\gtrsim M\log M$ zeros in the window $(M,2M]$, each carrying an
amplitude $a_\gamma\in[2/\sqrt5,\sqrt2\,]$ bounded away from $0$ and
$\infty$, force superexponential decay. Both $\sigma_m^2$ and $S_4(m)$
turn out to be exact values of the completed $L$-function
(Proposition~\ref{prop:varid} and Lemma~\ref{lem:S4mom}), so every
coefficient below is a closed form.

\subsection*{Setting and statement}

Let $\chi$ be a real primitive character mod $q\ge3$,
$\chi(-1)=(-1)^a$ with $a\in\{0,1\}$, and let
$\xi(s,\chi)=(q/\pi)^{(s+a)/2}\Gamma\bigl(\tfrac{s+a}2\bigr)L(s,\chi)$
be its completed $L$-function. We assume GRH$(\chi)$ and
$L(\tfrac12,\chi)\ne0$ (for $q=3,4$ the nonvanishing is unconditional,
as noted in the proof of Proposition~\ref{prop:varasymp}); $\gamma$
runs over the positive ordinates of $L(s,\chi)$, listed with
multiplicity. For $m>0$ put $M=m+\tfrac12$,
\[
a_\gamma(m)=\frac{2M}{\sqrt{M^2+\gamma^2}},\qquad
X_{m,\chi}=1+V_m,\qquad
V_m=\sum_{\gamma>0}2a_\gamma(m)\cos\theta_\gamma,
\]
with $\theta_\gamma$ i.i.d.\ uniform on $[0,2\pi)$ (the series
converges a.s.\ and in $L^2$, as in Section~\ref{sec:theorem}), and
\[
\delta_q(m)\;:=\;\mathbb P\bigl(X_{m,\chi}>0\bigr),
\qquad
\sigma_m^2\;:=\;\sum_{\gamma>0}2a_\gamma(m)^2,
\qquad
S_4(m)\;:=\;\sum_{\gamma>0}\bigl(2a_\gamma(m)\bigr)^4 .
\]
For $q=4$, Theorem~\ref{thm:interp}(ii) identifies $\delta_4(m)$, under
LI$(\chi_4)$, with the logarithmic density $\delta(m)$ of the set where
the weighted race $R_m$ is positive; for other moduli $\delta_q(m)$ is,
in this paper, a statement about the limiting random variable only (see
the Scope remark at the end of the section). Note that the theorem
below does \emph{not} assume LI: independence of the phases is built
into the definition of $X_{m,\chi}$, and LI is what transfers the
conclusion to the race.

\begin{theorem}[dissolution law]\label{thm:dissolution}
Let $\chi$ be a real primitive character mod $q$; assume GRH$(\chi)$
and $L(\tfrac12,\chi)\ne0$. Then, with
$\sigma_m^2=4M\,\frac{\xi'}{\xi}(m+1,\chi)$ as in
\eqref{eq:varidq} and $S_4(m)$ as in Lemma~\ref{lem:S4mom}, as
$m\to\infty$:
\begin{itemize}
\item[(i)]
$\displaystyle
\delta_q(m)-\tfrac12\;=\;\frac{1}{\sqrt{2\pi\sigma_m^2}}
\Bigl(1+O\Bigl(\frac{1}{M\log M}\Bigr)\Bigr);
$
\item[(ii)] more precisely,
$\displaystyle
\delta_q(m)-\tfrac12\;=\;\frac{1}{\sqrt{2\pi\sigma_m^2}}
\Bigl(1-\frac{1}{6\sigma_m^2}-\frac{3S_4(m)}{64\,\sigma_m^4}
+O\Bigl(\frac{1}{(M\log M)^2}\Bigr)\Bigr);
$
\item[(iii)] consequently, by
Proposition~\ref{prop:varasymp},
\[
\delta_q(m)-\tfrac12\;=\;
\Bigl(4\pi M\log\tfrac{qM}{2\pi}\Bigr)^{-1/2}
\Bigl(1+O\Bigl(\frac{1}{M\log M}\Bigr)\Bigr).
\]
\end{itemize}
All implied constants depend on $q$ alone; they are effectively
computable in principle and we have not optimized them. Nothing is
claimed uniformly in $q$.
\end{theorem}

\begin{corollary}[the race itself, $q=4$]\label{cor:dissolution-race}
Assume GRH$(\chi_4)$ and LI$(\chi_4)$. Then the logarithmic density
$\delta(m)$ of $\{x\ge2:R_m(x)>0\}$ of Theorem~\ref{thm:interp}
satisfies (i)--(iii) of Theorem~\ref{thm:dissolution} verbatim. In
particular $\delta(m)\to\tfrac12$ as $m\to\infty$ at the rate
$(4\pi M\log\tfrac{2M}{\pi})^{-1/2}$, closing the gap left open in
Theorem~\ref{thm:interp}(iii).
\end{corollary}

\begin{corollary}[closed second order]\label{cor:secondorder}
Under the hypotheses of Theorem~\ref{thm:dissolution}, with
$L:=\log\tfrac{qM}{2\pi}$,
\[
\delta_q(m)-\tfrac12\;=\;\frac{1}{\sqrt{2\pi\sigma_m^2}}
\Bigl(1-\frac{11}{24\,ML}+O\Bigl(\frac{1}{ML^2}\Bigr)\Bigr)
\qquad(m\to\infty).
\]
\end{corollary}

The relative error in (i) is thus not just $O(1/(M\log M))$ but
asymptotically equal to $-\tfrac{11}{24}(ML)^{-1}$: the density
approaches the Gaussian prediction from \emph{below}, as every entry
of Table~\ref{tab:dissolution} shows.

We assemble the proof from four lemmas: elementary Bessel estimates
(Lemma~\ref{lem:J0facts}); a count of zeros in the dyadic window
$(M,2M]$ (Lemma~\ref{lem:window}); the closed form and asymptotics of
$S_4$ (Lemma~\ref{lem:S4mom}); and the inversion formula
(Lemma~\ref{lem:invert}) together with the tail bound that justifies
it (Lemma~\ref{lem:bigt}).

\subsection*{Lemmas}

\begin{lemma}[Bessel facts]\label{lem:J0facts}
Let $J_0(z)=\sum_{k\ge0}(-1)^k(z^2/4)^k/(k!)^2$. For real $z$:
\begin{itemize}
\item[(J1)] $|J_0(z)|\le1$, with strict inequality for $z\ne0$;
moreover for every $z_0>0$,
\[
\rho(z_0)\;:=\;\sup_{z\ge z_0}|J_0(z)|\;<\;1 .
\]
\item[(J2)] $|J_0(z)|\le\dfrac{8}{\pi}\,z^{-1/2}$ for all $z>0$.
\item[(J3)] For $0\le z\le1$: $J_0(z)\ge\tfrac34$ and
$\bigl|\log J_0(z)+\tfrac{z^2}4\bigr|\le\tfrac{z^4}{17}$.
\item[(J4)] For $0\le z\le1$:
$\bigl|\log J_0(z)+\tfrac{z^2}4+\tfrac{z^4}{64}\bigr|
\le\tfrac{z^6}{80}$.
\end{itemize}
\end{lemma}

\begin{proof}
Throughout we use the integral representation
\begin{equation}\label{eq:J0int}
J_0(z)\;=\;\frac2\pi\int_0^{\pi/2}\cos(z\sin\varphi)\,d\varphi
\;=\;\frac2\pi\int_0^1\frac{\cos(zu)}{\sqrt{1-u^2}}\,du ,
\end{equation}
obtained by expanding $\cos(z\sin\varphi)$ into its (uniformly
convergent) power series and integrating termwise with Wallis's
formula $\int_0^{\pi/2}\sin^{2k}\varphi\,d\varphi
=\tfrac\pi2\binom{2k}{k}4^{-k}$, which reproduces the defining series;
the second form is the substitution $u=\sin\varphi$.

(J1) The first form of \eqref{eq:J0int} gives $|J_0(z)|\le1$. If
$z\ne0$, then $\varphi\mapsto z\sin\varphi$ is injective on
$[0,\pi/2]$, so $\cos(z\sin\varphi)=\pm1$ for at most finitely many
$\varphi$, and by continuity $|J_0(z)|<1$ strictly. For the supremum:
by (J2), $|J_0(z)|\le\tfrac12$ for $z\ge26>(16/\pi)^2$; on the compact
interval $[z_0,26]$ the continuous function $|J_0|$ attains a maximum,
which is $<1$ by the strict inequality just proved. Hence
$\rho(z_0)=\max\bigl(\max_{[z_0,26]}|J_0|,\tfrac12\bigr)<1$ (for
$z_0\ge26$ the compact piece is vacuous and $\rho(z_0)\le\tfrac12$
directly). (No
numerical value of $\rho$ is used in the proofs; for orientation,
$\rho(1/(2\sqrt5))=J_0(1/(2\sqrt5))=0.98754$, computed in the
verification script.)

(J2) For $0<z<1$ the bound exceeds $8/\pi>1$ and is trivial. For
$z\ge1$ set $\delta=1/z\in(0,1]$ and split \eqref{eq:J0int} at
$u=1-\delta$. The singular piece:
\[
\int_{1-\delta}^1\frac{du}{\sqrt{1-u^2}}
\le\int_{1-\delta}^1\frac{du}{\sqrt{1-u}}=2\sqrt\delta .
\]
The oscillatory piece, integrating by parts
($\cos(zu)=\tfrac1z\frac{d}{du}\sin(zu)$) and using
$1-u^2\ge\delta(2-\delta)\ge\delta$ at $u=1-\delta$ together with
$\int_0^{1-\delta}u(1-u^2)^{-3/2}du=(1-u^2)^{-1/2}\big|_0^{1-\delta}
\le\delta^{-1/2}$:
\[
\Bigl|\int_0^{1-\delta}\frac{\cos(zu)}{\sqrt{1-u^2}}\,du\Bigr|
\;\le\;\frac1z\cdot\frac1{\sqrt\delta}
+\frac1z\int_0^{1-\delta}\frac{u\,du}{(1-u^2)^{3/2}}
\;\le\;\frac2{z\sqrt\delta} .
\]
With $\delta=1/z$ both pieces are $2z^{-1/2}$, so
$|J_0(z)|\le\tfrac2\pi\cdot4z^{-1/2}=\tfrac8\pi z^{-1/2}$.

(J3) For $|z|\le1$ the series terms $c_k=(z^2/4)^k/(k!)^2$ decrease
strictly ($c_{k+1}/c_k=z^2/(4(k+1)^2)\le\tfrac14$), so the alternating
series gives
\[
1-\tfrac{z^2}4\;\le\;J_0(z)\;\le\;1-\tfrac{z^2}4+\tfrac{z^4}{64},
\qquad\text{hence } J_0(z)\ge\tfrac34,
\]
and, writing $J_0=1-w$ with $w=\tfrac{z^2}4-R$,
\[
R:=\sum_{k\ge2}(-1)^kc_k\in\bigl[0,\tfrac{z^4}{64}\bigr],
\qquad
w\in\bigl[\tfrac{z^2}4-\tfrac{z^4}{64},\,\tfrac{z^2}4\bigr]
\subset\bigl[0,\tfrac14\bigr].
\]
Then $\log J_0+\tfrac{z^2}4=R-\sum_{j\ge2}\tfrac{w^j}j$, and
$\sum_{j\ge2}\tfrac{w^j}j\le\tfrac{w^2}{2(1-w)}\le\tfrac23w^2
\le\tfrac23\bigl(\tfrac{z^2}4\bigr)^2=\tfrac{z^4}{24}$, so the total
is at most $\tfrac{z^4}{64}+\tfrac{z^4}{24}=\tfrac{11z^4}{192}
\le\tfrac{z^4}{17}$.

(J4) Refine each piece by one order. First,
$R-\tfrac{z^4}{64}=\sum_{k\ge3}(-1)^kc_k$ has absolute value at most
$c_3=\tfrac{z^6}{2304}$. Next,
$w^2=\tfrac{z^4}{16}-\tfrac{z^2R}2+R^2$, so
\[
\log J_0+\tfrac{z^2}4+\tfrac{z^4}{64}
=\Bigl(R-\tfrac{z^4}{64}\Bigr)+\tfrac{z^2R}4-\tfrac{R^2}2
-\sum_{j\ge3}\tfrac{w^j}j ,
\]
and, for $|z|\le1$: $\tfrac{z^2R}4\le\tfrac{z^6}{256}$,
$\tfrac{R^2}2\le\tfrac{z^8}{8192}\le\tfrac{z^6}{8192}$, and
$\sum_{j\ge3}\tfrac{w^j}j\le\tfrac{w^3}{3(1-w)}
\le\tfrac49\bigl(\tfrac{z^2}4\bigr)^3=\tfrac{z^6}{144}$. The sum of
the four bounds is
$z^6\bigl(\tfrac1{2304}+\tfrac1{256}+\tfrac1{8192}+\tfrac1{144}\bigr)
<0.0115\,z^6<\tfrac{z^6}{80}$. (The true limit ratio as $z\to0$ is
$\tfrac1{576}$; the script checks the claimed bounds on a grid in
multiprecision, float64 being insufficient near $z=0$.)
\end{proof}

\begin{lemma}[zeros in the dyadic window]\label{lem:window}
Unconditionally, with
$N_\chi(T):=\#\{\gamma\in(0,T]\}$, there is $M_0(q)$ such that for all
$M\ge M_0(q)$,
\[
N_0\;:=\;N_\chi(2M)-N_\chi(M)\;\ge\;\frac{M}{4\pi}\log M
\;\ (\ge5).
\]
\end{lemma}

\begin{proof}
For a primitive $\chi$ mod $q$ the Riemann--von Mangoldt formula gives
$N_\chi(T)=\frac{T}{2\pi}\log\frac{qT}{2\pi e}+O(\log qT)$ for
$T\ge2$ \cite[\S16]{Dav00} (Davenport counts zeros with $|\gamma|\le
T$; for real $\chi$ the multiset of ordinates is symmetric under
$\gamma\mapsto-\gamma$, as used in Proposition~\ref{prop:varid}, so
the one-sided count is half). Hence
\[
N_\chi(2M)-N_\chi(M)
=\frac{M}{2\pi}\Bigl(\log\frac{qM}{2\pi e}+2\log2\Bigr)+O_q(\log M)
\ge\frac{M}{2\pi}\log\frac{qM}{2\pi e}-C_q\log M ,
\]
which exceeds $\frac{M}{4\pi}\log M$ once $M\ge M_0(q)$.
\end{proof}

\begin{lemma}[fourth moment of the amplitudes]\label{lem:S4mom}
Assume GRH$(\chi)$ and $L(\tfrac12,\chi)\ne0$. Then for every
$m>-\tfrac12$,
\begin{equation}\label{eq:S4id}
S_4(m)\;=\;\sum_{\gamma>0}\frac{256\,M^4}{(M^2+\gamma^2)^2}
\;=\;16\,\sigma_m^2\;-\;64\,M^2\,(\log\xi)''(m+1,\chi),
\end{equation}
and, as $m\to\infty$,
\begin{equation}\label{eq:S4asymp}
S_4(m)\;=\;32M\Bigl(\log\frac{qM}{2\pi}-1\Bigr)+32(2a-1)
+O_q\Bigl(\frac1M\Bigr).
\end{equation}
Moreover, unconditionally on the size of $m$,
$S_4(m)\le32\,\sigma_m^2$ and
$S_6(m):=\sum_\gamma(2a_\gamma)^6\le16\,S_4(m)$.
\end{lemma}

\begin{proof}
The trivial bounds are immediate from $(2a_\gamma)^2\le16$ and
$\sum_\gamma(2a_\gamma)^2=2\sigma_m^2$.

For \eqref{eq:S4id}: exactly as in the proof of the closed form for
$C$ (Corollary~\ref{cor:Cexact}, applied to $\chi$ in place of
$\chi_4$ as in Proposition~\ref{prop:varasymp}), the series
$(\log\Xi)'(x)=\sum_{\gamma>0}2x/(x^2+\gamma^2)$,
$\Xi(z):=\xi(\tfrac12+z,\chi)$, may be differentiated termwise on a
neighbourhood of $x=M$: the differentiated terms
$2(\gamma^2-x^2)/(x^2+\gamma^2)^2$ are dominated by
$2/(x^2+\gamma^2)\le2/\gamma^2$, uniformly on compacts. Hence
\[
(\log\Xi)''(M)=\sum_{\gamma>0}\frac{2(\gamma^2-M^2)}{(M^2+\gamma^2)^2},
\]
and the identity is termwise algebra:
\[
16\cdot\frac{8M^2}{M^2+\gamma^2}
-64M^2\cdot\frac{2(\gamma^2-M^2)}{(M^2+\gamma^2)^2}
=\frac{128M^2(M^2+\gamma^2)-128M^2\gamma^2+128M^4}{(M^2+\gamma^2)^2}
=\frac{256M^4}{(M^2+\gamma^2)^2},
\]
summed over $\gamma>0$, using
$\sigma_m^2=\sum_\gamma8M^2/(M^2+\gamma^2)$
(Proposition~\ref{prop:varid}, \eqref{eq:varidq}) and
$(\log\Xi)''(M)=(\log\xi)''(m+1,\chi)$.

For \eqref{eq:S4asymp}: $(\log\xi)''(s)=\tfrac14\psi'\bigl(\tfrac{s+a}2
\bigr)+\bigl(\tfrac{L'}{L}\bigr)'(s,\chi)$. At $s=m+1$,
$x:=\tfrac{m+1+a}2=\tfrac M2+\tfrac{2a+1}4$, and Euler--Maclaurin for
$\psi'(x)=\sum_{n\ge0}(n+x)^{-2}$ (comparison with
$\int_0^\infty(t+x)^{-2}dt=1/x$; correction $f(0)/2=1/(2x^2)$;
remainder at most
$\tfrac1{12}|f'(0)|+\tfrac1{12}\int_0^\infty|f''|\le\tfrac1{2x^3}$)
gives $\psi'(x)=\tfrac1x+\tfrac1{2x^2}+O(x^{-3})$, whence
\[
\tfrac14\psi'(x)=\frac1{2M}\cdot\frac1{1+\frac{2a+1}{2M}}
+\frac1{2M^2}+O\Bigl(\frac1{M^3}\Bigr)
=\frac1{2M}-\frac{2a+1}{4M^2}+\frac1{2M^2}
+O\Bigl(\frac1{M^3}\Bigr)
=\frac1{2M}+\frac{1-2a}{4M^2}+O\Bigl(\frac1{M^3}\Bigr),
\]
so $64M^2\cdot\tfrac14\psi'(x)=32M+16(1-2a)+O(1/M)$. The Dirichlet
series $(\tfrac{L'}{L})'(s,\chi)=\sum_{n\ge2}\Lambda(n)\chi(n)(\log
n)\,n^{-s}$ is bounded at $s=m+1$ by
$2^{-(m-1)}\sum_n\Lambda(n)(\log n)n^{-2}\ll2^{-m}$, so
$64M^2(\tfrac{L'}{L})'(m+1)\ll M^22^{-m}=O(1/M)$ eventually. Combining
with $16\sigma_m^2=32M\log\tfrac{qM}{2\pi}+16(2a-1)+O(1/M)$
(Proposition~\ref{prop:varasymp}) yields \eqref{eq:S4asymp}:
\[
S_4(m)=32M\log\tfrac{qM}{2\pi}+16(2a-1)-32M-16(1-2a)+O(1/M).
\qedhere
\]
\end{proof}

Numerically \eqref{eq:S4id} and \eqref{eq:S4asymp} agree to the
expected order: at $(q,m)=(4,100)$, the closed form gives
$S_4=10190.02$, the $511$ ordinates plus Riemann--von Mangoldt tail
give $10190.11$, and \eqref{eq:S4asymp} without its $O(1/M)$ gives
$10189.98$.

\begin{lemma}[characteristic function and inversion]\label{lem:invert}
Let $V_m=\sum_{\gamma>0}2a_\gamma(m)\cos\theta_\gamma$ as above. Then:
\begin{itemize}
\item[(a)] the law of $V_m$ is symmetric, and its characteristic
function is
$\Phi_m(t)=\prod_{\gamma>0}J_0(2a_\gamma(m)t)$, the limit of the
partial products, which is real and even; moreover for every finite
set $F$ of positive ordinates,
$|\Phi_m(t)|\le\prod_{\gamma\in F}|J_0(2a_\gamma(m)t)|$;
\item[(b)] if $\Phi_m\in L^1(\mathbb R)$, then $V_m$ has a bounded
continuous even density and
\[
\delta_q(m)-\tfrac12\;=\;\mathbb P(V_m>-1)-\tfrac12
\;=\;\frac1\pi\int_0^\infty\Phi_m(t)\,\frac{\sin t}{t}\,dt ,
\]
the integral converging absolutely.
\end{itemize}
\end{lemma}

\begin{proof}
(a) For a single summand, $\mathbb E\,e^{it\cdot2a\cos\theta}
=\frac1{2\pi}\int_0^{2\pi}e^{2iat\cos\theta}d\theta=J_0(2at)$: expand
the exponential and integrate termwise (odd powers vanish;
$\frac1{2\pi}\int_0^{2\pi}\cos^{2k}\theta\,d\theta
=\binom{2k}{k}4^{-k}$), recovering the series of $J_0$; in particular
each summand is symmetric ($\theta\mapsto\theta+\pi$ preserves the
uniform law and flips the sign), hence so are the partial sums
$S_N=\sum_{n\le N}2a_{\gamma_n}\cos\theta_{\gamma_n}$ and their a.s.\
limit $V_m$. Independence gives the characteristic function of $S_N$
as the $N$-th partial product; $S_N\to V_m$ a.s.\ implies convergence
in distribution, so the characteristic functions converge pointwise
to that of $V_m$, which is therefore the limit of the partial
products, real (symmetry) and even. For the subset bound, take
$N$ large enough that the first $N$ ordinates contain $F$; all
factors have modulus $\le1$ by (J1), so the modulus of the $N$-th
partial product is at most $\prod_{\gamma\in F}|J_0|$, and let
$N\to\infty$.

(b) By the Fourier inversion theorem for integrable characteristic
functions \cite[Ch.~XV.3]{Fel71}, the law of $V_m$ has the bounded
continuous density $f(x)=\frac1{2\pi}\int_{\mathbb R}
e^{-itx}\Phi_m(t)\,dt$, which is even since $\Phi_m$ is real and
even. Hence, using symmetry and then Fubini (justified by
$\int_{\mathbb R}\int_{-1}^1|\Phi_m(t)|\,dx\,dt
=2\|\Phi_m\|_1<\infty$),
\[
\mathbb P(V_m>-1)=\tfrac12+\tfrac12\,\mathbb P(|V_m|<1)
=\tfrac12+\tfrac12\int_{-1}^1 f
=\tfrac12+\frac1{4\pi}\int_{\mathbb R}\Phi_m(t)\Bigl(\int_{-1}^1
e^{-itx}dx\Bigr)dt ,
\]
and $\int_{-1}^1e^{-itx}dx=2\sin t/t$, so
$\mathbb P(V_m>-1)-\tfrac12=\frac1{2\pi}\int_{\mathbb R}\Phi_m(t)
\frac{\sin t}t\,dt=\frac1\pi\int_0^\infty\Phi_m(t)\frac{\sin t}t\,dt$
by evenness. Finally $\delta_q(m)=\mathbb P(1+V_m>0)
=\mathbb P(V_m>-1)$.
\end{proof}

\begin{lemma}[superexponential decay for $t\ge\tfrac18$]\label{lem:bigt}
Assume GRH$(\chi)$. Let $\rho:=\rho\bigl(1/(2\sqrt5)\bigr)<1$ and
$\tilde\rho:=\max(\rho,\tfrac12)<1$ (absolute constants, by (J1)).
There is $M_0(q)$ such that for all $m$ with $M\ge M_0(q)$,
\[
\int_{1/8}^\infty|\Phi_m(t)|\,\frac{|\sin t|}{t}\,dt
\;\le\;17\,\tilde\rho^{\,N_0}
\;\le\;\exp\Bigl(-\frac{\log(1/\tilde\rho)}{8\pi}\,M\log M\Bigr),
\]
with $N_0=N_\chi(2M)-N_\chi(M)$ as in Lemma~\ref{lem:window};
in particular $\Phi_m\in L^1(\mathbb R)$.
\end{lemma}

\begin{proof}
For $\gamma\in(M,2M]$ we have $M^2+\gamma^2\in(2M^2,5M^2]$, so
\[
a_\gamma(m)=\frac{2M}{\sqrt{M^2+\gamma^2}}\in
\Bigl[\frac{2}{\sqrt5},\sqrt2\Bigr),
\qquad
2a_\gamma(m)\,t\;\ge\;\frac{4}{\sqrt5}\,t .
\]
By Lemma~\ref{lem:invert}(a) with $F=\{\gamma\in(M,2M]\}$ (a set of
$N_0$ ordinates, nonempty for $M\ge M_0(q)$):
\emph{(i)} for $t\ge\tfrac18$, every factor has argument
$\ge\tfrac{4}{\sqrt5}\cdot\tfrac18=\tfrac1{2\sqrt5}$, so by (J1)
$|\Phi_m(t)|\le\rho^{N_0}$;
\emph{(ii)} for any $t>0$, by (J2),
\[
|\Phi_m(t)|\;\le\;
\Bigl(\frac{8}{\pi}\Bigl(\frac{4t}{\sqrt5}\Bigr)^{-1/2}\Bigr)^{N_0}
=\Bigl(\frac{c_3}{\sqrt t}\Bigr)^{N_0},
\qquad
c_3:=\frac{4\cdot5^{1/4}}{\pi}=1.9040\ldots<2 ,
\]
so that $c_3/\sqrt t\le\tfrac12$ for $t\ge16$. Using
$|\sin t|/t\le1$ on $[\tfrac18,16]$ and $|\sin t|\le1$ beyond, and
(for $M\ge M_0(q)$, by Lemma~\ref{lem:window}) $N_0\ge5$:
\[
\int_{1/8}^{16}|\Phi_m|\,\frac{|\sin t|}t\,dt\le16\rho^{N_0},
\qquad
\int_{16}^\infty|\Phi_m|\,\frac{dt}t
\le2^{5-N_0}c_3^5\int_{16}^\infty t^{-7/2}\,dt
=2^{5-N_0}c_3^5\cdot\frac{2}{5\cdot1024}
\le\frac{2^{-N_0}}{2},
\]
where the middle step keeps five factors $(c_3/\sqrt t)^5$ and bounds
the remaining $N_0-5$ by $2^{-(N_0-5)}$. The sum is at most
$17\tilde\rho^{\,N_0}$, and with
$N_0\ge\tfrac{M}{4\pi}\log M$ the final bound follows for
$M\ge M_0(q)$ enlarged if necessary (to absorb the factor $17$ into
half the exponent). The same estimates with $|\sin t|/t$ replaced by
$1$ on $[\tfrac18,16]$ and by nothing beyond show
$\int_{1/8}^\infty|\Phi_m|<\infty$; on $[0,\tfrac18]$, $|\Phi_m|\le1$.
Hence $\Phi_m\in L^1$.
\end{proof}

\subsection*{Proof of Theorem~\ref{thm:dissolution}}

\begin{proof}
Throughout, $M\ge M_0(q)$ so that Lemmas~\ref{lem:window} and
\ref{lem:bigt} apply, and $G:=(2\pi\sigma_m^2)^{-1/2}$. By
Lemma~\ref{lem:bigt}, $\Phi_m\in L^1$, so Lemma~\ref{lem:invert}(b)
gives
\begin{equation}\label{eq:GP}
\delta_q(m)-\tfrac12
=\frac1\pi\int_0^{1/8}\Phi_m(t)\frac{\sin t}t\,dt
+O\bigl(e^{-c_qM\log M}\bigr),
\qquad c_q:=\frac{\log(1/\tilde\rho)}{8\pi}.
\end{equation}

\emph{The Gaussian factorization on $[0,\tfrac18]$.} For
$0\le t\le\tfrac18$ every argument satisfies
$2a_\gamma(m)t\le4t\le\tfrac12$, since $a_\gamma\le2$. By (J3) each
factor is $\ge\tfrac34>0$ and
$\sum_\gamma|\log J_0(2a_\gamma t)|\le\sum_\gamma\bigl(a_\gamma^2t^2
+(2a_\gamma t)^4/17\bigr)<\infty$, so the limit of the partial
products is
\begin{equation}\label{eq:smallt}
\Phi_m(t)=\exp\Bigl(-\frac{\sigma_m^2t^2}{2}+E(t)\Bigr),
\qquad
E(t):=\sum_{\gamma>0}\Bigl(\log J_0(2a_\gamma t)
+\frac{(2a_\gamma t)^2}{4}\Bigr),
\end{equation}
using $\sum_\gamma(2a_\gamma)^2t^2/4=\sigma_m^2t^2/2$. By (J3),
$|E(t)|\le\tfrac{t^4}{17}S_4$, and since $S_4\le32\sigma_m^2$
(Lemma~\ref{lem:S4mom}) and $t\le\tfrac18$,
\begin{equation}\label{eq:Eunif}
|E(t)|\;\le\;\frac{t^2}{17}\cdot\frac{1}{64}\cdot32\,\sigma_m^2
\;=\;\frac{\sigma_m^2t^2}{34}
\qquad\bigl(0\le t\le\tfrac18\bigr),
\end{equation}
so that $|\Phi_m(t)|\le e^{-(8/17)\sigma_m^2t^2}$ there: the Gaussian
decay survives the correction uniformly.

\emph{Decomposition.} Write, on $[0,\tfrac18]$,
$\Phi_m\frac{\sin t}t
=e^{-\sigma^2t^2/2}
-e^{-\sigma^2t^2/2}\bigl(1-\tfrac{\sin t}t\bigr)
+e^{-\sigma^2t^2/2}\bigl(e^{E}-1\bigr)\tfrac{\sin t}t$
(abbreviating $\sigma=\sigma_m$), so that
\begin{equation}\label{eq:decomp}
\frac1\pi\int_0^{1/8}\Phi_m\frac{\sin t}t\,dt
\;=\;G-T_1-T_2+T_3 ,
\end{equation}
with, using $\frac1\pi\int_0^\infty e^{-\sigma^2t^2/2}dt=G$ and the
moments $\frac1\pi\int_0^\infty t^{2k}e^{-\sigma^2t^2/2}dt
=(2k-1)!!\,G\,\sigma^{-2k}$:
\begin{itemize}
\item[$T_1$:] the completed Gaussian tail,
$0\le T_1=\frac1\pi\int_{1/8}^\infty e^{-\sigma^2t^2/2}dt
\le\frac{8}{\pi\sigma^2}e^{-\sigma^2/128}$;
\item[$T_2$:] the $\sin$ correction: since
$0\le1-\frac{\sin t}t\le\frac{t^2}6$ for all $t>0$ (from
$t-\frac{t^3}6\le\sin t\le t$),
\[
0\;\le\;T_2\;=\;\frac1\pi\int_0^{1/8}
e^{-\sigma^2t^2/2}\Bigl(1-\frac{\sin t}t\Bigr)dt
\;\le\;\frac1{6\pi}\int_0^{\infty}t^2e^{-\sigma^2t^2/2}\,dt
\;=\;\frac{G}{6\sigma^{2}} ;
\]
\item[$T_3$:] the Bessel correction: by \eqref{eq:Eunif} and
$|e^w-1|\le|w|e^{|w|}$,
\[
|T_3|\le\frac1\pi\int_0^{1/8}e^{-\sigma^2t^2/2}
\cdot\frac{t^4S_4}{17}\,e^{\sigma^2t^2/34}\,dt
\le\frac{S_4}{17\pi}\int_0^\infty t^4e^{-\frac{8}{17}\sigma^2t^2}dt
=\frac{3S_4}{8\cdot17\sqrt\pi}
\Bigl(\frac{17}{8}\Bigr)^{5/2}\sigma^{-5}
\le\frac{S_4}{12\,\sigma^{5}}\le\frac{G\,S_4}{4\,\sigma^{4}} ,
\]
the last step by $\sigma^{-5}=\sqrt{2\pi}\,G\sigma^{-4}$ and
$\sqrt{2\pi}/12<\tfrac14$.
\end{itemize}

\emph{Part (i).} By Proposition~\ref{prop:varasymp} there is
$M_1(q)\ge M_0(q)$ with
$ML\le\sigma_m^2\le3ML$ for $M\ge M_1(q)$, where
$L=\log\tfrac{qM}{2\pi}\asymp_q\log M$; and
$S_4\le33ML$ by \eqref{eq:S4asymp}. Hence
\[
\frac{T_2}{G}\le\frac1{6ML},\qquad
\frac{|T_3|}{G}\le\frac{33ML}{4M^2L^2}\le\frac{9}{ML},\qquad
\frac{T_1+O(e^{-c_qM\log M})}{G}
\ll\sigma\,e^{-\sigma^2/128}+\sigma\,e^{-c_qM\log M}
\le\frac1{ML}
\]
for $M$ large, and \eqref{eq:GP}--\eqref{eq:decomp} give
$\delta_q(m)-\tfrac12=G\bigl(1+O_q(1/(ML))\bigr)$, which is (i).

\emph{Part (ii).} Refine $T_2$ and $T_3$ by one order each. For
$T_2$: on $[0,\tfrac18]$ the alternating series of $\frac{\sin t}t$
has decreasing terms, so
$\bigl|1-\frac{\sin t}t-\frac{t^2}6\bigr|\le\frac{t^4}{120}$, and
extending the integral to $\infty$ costs $O(e^{-\sigma^2/129})$:
\[
T_2=\frac1{6\pi}\int_0^{\infty}t^2e^{-\sigma^2t^2/2}dt
+O\Bigl(\frac1{120\pi}\int_0^\infty t^4e^{-\sigma^2t^2/2}dt\Bigr)
+O\bigl(e^{-\sigma^2/129}\bigr)
=\frac{G}{6\sigma^2}+O\bigl(G\sigma^{-4}\bigr).
\]
For $T_3$: by (J4), $E(t)=-\frac{t^4}{64}S_4+E_2(t)$ with
$|E_2(t)|\le\frac{t^6}{80}S_6\le\frac{t^6}{5}S_4$
(Lemma~\ref{lem:S4mom}); and $e^E-1=E+\omega$ with
$|\omega|\le\frac{E^2}2e^{|E|}
\le\frac{t^8S_4^2}{578}\,e^{\sigma^2t^2/34}$ by \eqref{eq:Eunif}.
Hence
\[
T_3=-\frac{S_4}{64\pi}\int_0^{1/8}t^4e^{-\sigma^2t^2/2}
\frac{\sin t}t\,dt
+O\Bigl(\frac{S_4}{5\pi}\int_0^\infty
t^6e^{-\sigma^2t^2/2}dt\Bigr)
+O\Bigl(\frac{S_4^2}{578\pi}\int_0^\infty
t^8e^{-\frac{8}{17}\sigma^2t^2}dt\Bigr),
\]
and in the first integral $\frac{\sin t}t=1+O(t^2)$, so, with the
moment values and extension to $\infty$ as before,
\[
T_3=-\frac{3\,S_4}{64}\,G\,\sigma^{-4}
+O\bigl(G\,S_4\,\sigma^{-6}\bigr)
+O\bigl(G\,S_4^2\,\sigma^{-8}\bigr).
\]
Since $\sigma^2\asymp_qML$ and $S_4\asymp_qML$, the error terms
$G\cdot O(\sigma^{-4}+S_4\sigma^{-6}+S_4^2\sigma^{-8})$ are all
$G\cdot O_q\bigl((ML)^{-2}\bigr)$, and the exponentially small terms
$T_1$, \eqref{eq:GP} are absorbed. Collecting,
\[
\delta_q(m)-\tfrac12
=G\Bigl(1-\frac1{6\sigma_m^2}-\frac{3S_4}{64\sigma_m^4}
+O_q\bigl((M\log M)^{-2}\bigr)\Bigr).
\]

\emph{Part (iii).} By Proposition~\ref{prop:varasymp},
$\sigma_m^2=2ML+O(1)$, so
$(2\pi\sigma_m^2)^{-1/2}=(4\pi ML)^{-1/2}(1+O(1/(ML)))$; combine with
(i).
\end{proof}

\begin{proof}[Proof of Corollary~\ref{cor:dissolution-race}]
$L(\tfrac12,\chi_4)>0$ unconditionally, so
Theorem~\ref{thm:dissolution} applies to $\chi_4$; and under
GRH$(\chi_4)$+LI$(\chi_4)$, Theorem~\ref{thm:interp}(ii) gives
$\delta(m)=\nu_m\bigl((0,\infty)\bigr)=\mathbb P(X_{m,\chi_4}>0)
=\delta_4(m)$.
\end{proof}

\begin{proof}[Proof of Corollary~\ref{cor:secondorder}]
Insert $\sigma_m^2=2ML+O(1)$ and
$S_4=32M(L-1)+O(1)$ (Proposition~\ref{prop:varasymp},
Lemma~\ref{lem:S4mom}) into (ii):
\[
\frac1{6\sigma_m^2}=\frac1{12ML}+O\Bigl(\frac1{M^2L^2}\Bigr),
\qquad
\frac{3S_4}{64\sigma_m^4}
=\frac{96M(L-1)}{256M^2L^2}+O\Bigl(\frac1{M^2L^2}\Bigr)
=\frac{3}{8ML}+O\Bigl(\frac1{ML^2}\Bigr),
\]
and $\frac1{12}+\frac38=\frac{11}{24}$; the error
$O((ML)^{-2})$ of (ii) is also $O(1/(ML^2))$.
\end{proof}

\begin{remark}[numerical consistency]\label{rem:dissolution-numerics}
Table~\ref{tab:dissolution} compares the theorem with deterministic
Gil--Pelaez reference densities (computational, not part of any
proof): $511$ resp.\ $537$ ordinates, unseen zeros Gaussianized with
the \emph{exact} tail variance of Remark~\ref{rem:tails}, trapezoidal
inversion with grid-halving and window-doubling stable to
$10^{-9}$. Both $\sigma_m^2$ and $S_4$ are evaluated by their closed
forms \eqref{eq:varidq}, \eqref{eq:S4id} --- no zero data enters the
predictions. Gaussianizing the unseen zeros biases the \emph{reference}
values themselves: on $t\le\tfrac18$ the $T_3$ estimate, applied to
the tail factors alone, rigorously bounds the relative bias by
$0.21\,S_4^{\mathrm{tail}}/\sigma_m^4$; on $t>\tfrac18$ that estimate
does not apply, and we instead bound the remaining bias by the same
integrand with the \emph{computed} $|\Phi_m|$ inserted, which is below
$9\times10^{-7}$ at every table point. The total ---
$9.9\times10^{-6}$ at $(q,m)=(4,100)$, less elsewhere in the table ---
is the accuracy floor of the comparison; the $t>\tfrac18$ piece is a
numerical evaluation, not a theorem, on the same footing as the
reference densities themselves. The certified enclosures of Table~\ref{tab:enclose},
computed by an independent route, contain the (unrounded) reference
values used here at all four table points (e.g.\
$0.51377836\in[0.5137782134,0.5137783837]$ at $(4,100)$; the table
entry $0.5137784$ is that reference rounded to seven decimals and lies
just above the enclosure's upper end, consistent with the one-sided
Gaussianization bias of the reference noted above).

\begin{table}[h]\centering
\begin{tabular}{lcccccc}
\toprule
$(q,m)$ & $\sigma_m^2$ & $S_4$ & $\delta_q(m)$ (ref.) &
$\dfrac{\delta_q(m)-\frac12}{(2\pi\sigma_m^2)^{-1/2}}-1$ &
$-\dfrac1{6\sigma_m^2}-\dfrac{3S_4}{64\sigma_m^4}$ &
residual \\
\midrule
$(4,8)$   & $29.7125$  & $219.83$   & $0.5718701$ & $-1.80\times10^{-2}$ & $-1.73\times10^{-2}$ & $-7.4\times10^{-4}$ \\
$(4,20)$  & $106.3260$ & $1061.32$  & $0.5384553$ & $-6.05\times10^{-3}$ & $-5.97\times10^{-3}$ & $-7.9\times10^{-5}$ \\
$(4,100)$ & $836.8744$ & $10190.02$ & $0.5137784$ & $-8.80\times10^{-4}$ & $-8.81\times10^{-4}$ & $+7.0\times10^{-7}$ \\
$(3,60)$  & $407.9701$ & $4607.56$  & $0.5197175$ & $-1.71\times10^{-3}$ & $-1.71\times10^{-3}$ & $-4.9\times10^{-6}$ \\
\bottomrule
\end{tabular}
\caption{The dissolution law against deterministic reference
densities. Column 5 is the observed relative deviation from the
leading term of Theorem~\ref{thm:dissolution}(i); column 6 is the
second-order term of (ii); the residual --- the relative error of the
two-term prediction, equal to
$(\text{col.~5}-\text{col.~6})/(1+\text{col.~6})$ --- shrinks like
$(M\log M)^{-2}$, as (ii) requires --- indeed at $(4,100)$ it is below
the $9.9\times10^{-6}$ modelling floor of the reference value, i.e.\
(ii) is confirmed there to the full accuracy of the computation. The
$m=100$ row is the consistency check promised in the introduction:
$(2\pi\cdot836.8744)^{-1/2}=0.0137905$ against
$\delta_4(100)-\tfrac12=0.0137784$, a relative gap of
$8.8\times10^{-4}$, against the theorem's error scale
$1/(M\log\tfrac{qM}{2\pi})=1/418=2.4\times10^{-3}$.}
\label{tab:dissolution}
\end{table}

Corollary~\ref{cor:secondorder} at $(4,100)$: $-\tfrac{11}{24}(ML)^{-1}
=-1.10\times10^{-3}$ against the observed $-8.80\times10^{-4}$; the
difference, $2.2\times10^{-4}$, is consistent with the corollary's
$O(1/(ML^2))=O(5.8\times10^{-4})$ --- at these heights $L\approx4.2$
is still small, which is why the table works with (ii) rather than
the corollary. The reference values $\delta_4(100)=0.5137784$ and
$\delta_3(60)=0.5197175$ refine the six-decimal values of the
introduction: $0.513778$ in Table~\ref{tab:delta}, and $0.519716$ in
the text beside it (the $(3,60)$ point is not a table entry; its value
shifts by $1.5\times10^{-6}$ upon replacing the estimated tail
variance of the original computation by the exact one).
\end{remark}

\begin{remark}[scope, and what is not proved]\label{rem:dissolution-scope}
(1) \emph{Race interpretation.} For $q=4$,
Corollary~\ref{cor:dissolution-race} makes the dissolution law a
statement about the weighted prime race itself. For $q=3$ the proof
of Theorem~\ref{thm:interp} goes through verbatim (nothing in
Lemmas~\ref{lem:trunc}--\ref{lem:squares} or the
\cite{ANS14} verification is specific to the modulus $4$ beyond
$\chi$ being real, primitive and nonprincipal), but we have not
restated it, so for $q=3$ --- and for every other real primitive
character --- Theorem~\ref{thm:dissolution} is, in this paper, a
theorem about the limiting random variable $X_{m,\chi}$ only. (2)
\emph{Hypotheses.} LI is not among the hypotheses of
Theorem~\ref{thm:dissolution}; GRH enters through the reality of the
ordinates and the closed forms \eqref{eq:varidq}, \eqref{eq:S4id},
and $L(\tfrac12,\chi)\ne0$ (automatic for $q=3,4$) through the
Hadamard product. (3) \emph{Constants.} The constants are effective
in principle ($\rho$ is a computable supremum, numerically $0.98754$;
the Riemann--von Mangoldt constants for fixed $q$ are classical) but
the large-$t$ bound uses only the zeros in $(M,2M]$ and is extremely
wasteful; no attempt has been made to make $M_0(q)$ small, and at the
heights of Table~\ref{tab:dissolution} the \emph{proved} bounds on
the tail terms are not yet numerically dominant --- the table is a
consistency check on the asymptotics, not a certified enclosure. (4)
\emph{Uniformity.} Nothing is claimed uniformly in $q$; a uniform
version would require the window count and
Proposition~\ref{prop:varasymp} with explicit $q$-dependence,
which we have not carried out (cf.\ the $q$-aspect asymptotics of
Fiorilli--Martin \cite{FM13}, whose densities-near-$\tfrac12$
regime is the natural companion statement). (5) The reference
densities of Remark~\ref{rem:dissolution-numerics} remain
conditional computations under GRH+LI, as throughout this
literature.
\end{remark}
